\documentclass[5pt]{article}
\usepackage{multicol,multirow}
\usepackage{graphicx} % Required for inserting images
\usepackage[margin=0.45cm]{geometry}
\usepackage{xcolor}
\usepackage{array}

\definecolor{LightGray}{gray}{0.9}
\usepackage[inline]{enumitem}
\usepackage{minted}

\newcolumntype{L}{>{\footnotesize}l}
\newcolumntype{C}{>{\ttfamily\footnotesize}l}

\begin{document}
\begin{center}
    \Large{\textbf{Rust Cheat Sheet}}\\
    \small{Last Updated: \today}\hfill\small{\textcopyright Maximilien Notz \the\year{}}
    \hspace*{-\dimexpr\oddsidemargin+1in\relax}\rule{\paperwidth}{0.4pt}
\end{center}


% content 
\begin{multicols}{2}
\setcounter{secnumdepth}{0}

\subsection{Cargo Commands}
Cargo is the default: package manager, builder for rust projects and version control via git. All of the following commend should be preced by \texttt{cargo}.\\
\begin{tabular}{C L}
    new new\_p  & Create a new project directory called \texttt{new\_p}.\\
    run         & Build and run the project.\\
    build       & Build the project.\\
    clean       & Remove the target directory.\\
    test        & Run the tests.\\
    update      & Update dependencies.\\
    doc         & Build documentation for the project.\\
\end{tabular}

\subsection{General Reminders}
\begin{tabular}{C L}
    use module::submodule & Import a module or submodule.\\
    println!("hello") & Print \texttt{hello} to the console.\\
\end{tabular}

\subsection{Strings}
\begin{tabular}{C L}
    String::new()     & Create a new empty string.\\
    String::from("s") & Create a new string with the value \texttt{"s"}.\\
    \{var\}           & Insert the value of \texttt{var} into the string.\\
\end{tabular}

\subsection{Input/Output}
For using input and output: the \texttt{std::io} module must be imported.\\
\begin{tabular}{C L}
    io::stdin()       & Get a handle to the standard input.\\
    io::stdout()      & Get a handle to the standard output.\\
\end{tabular}


\subsection{Basic syntax}
\subsubsection{Hello World}
\begin{minted}[bgcolor=LightGray]{rust}
fn main() {
    println!("Hello, world!");
}
\end{minted}

\begin{minted}[bgcolor=LightGray]{rust}
let mut s; // create a mutable variable s
let s1; // create a constant string s1
\end{minted}

\end{multicols}
\end{document}
